% Two loopback wires, a button, and the power stage that is not there.
\begin{tikzpicture}[font=\sffamily\small, x=1cm, y=1cm]

\node[board, minimum width=52mm, minimum height=30mm] (nucleo) at (0,0) {};
\node[text=white, font=\sffamily\bfseries\small] at (0,1.0) {NUCLEO-H7A3ZI-Q};
\node[text=white, font=\sffamily\tiny] at (0,0.62) {advanced timer: three pairs, dead time, break};
\node[pinrow, minimum width=46mm] at (0,1.38) {};

% the six outputs, as a row
\foreach \i/\lab in {0/{H1}, 1/{L1}, 2/{H2}, 3/{L2}, 4/{H3}, 5/{L3}}{
  \node[draw=white, fill=black!45, text=white, minimum width=6mm,
        minimum height=5mm, font=\sffamily\tiny] at (-1.5+\i*0.6,-0.1) {\lab};
}
\node[text=white, font=\sffamily\tiny] at (0,-0.65) {three complementary pairs};

% the loopback
\draw[cable, draw=red!60!black] (-1.5,-0.35) -- ++(0,-0.75) -- ++(3.9,0) -- ++(0,1.1);
\draw[cable, draw=blue!55!black] (-0.9,-0.35) -- ++(0,-1.05) -- ++(4.3,0) -- ++(0,1.4);
\node[text=white, font=\sffamily\tiny] at (0.2,-1.28) {loopback: H1 and L1 to two captures};

% the button as break input
\node[ledsym, fill=black!30] (btn) at (2.0,0.22) {};
\node[text=white, font=\sffamily\tiny, anchor=east] at (1.8,0.22) {PC13 as the break input};

% ---------------- what is not there
\node[absentblk, text width=38mm] (stage) at (6.6,1.9)
  {gate driver and half bridge, three of each\\{\scriptsize plus a supply, a shunt
   and an amplifier for current sensing}};
\node[absentblk, text width=38mm] (motor) at (6.6,-0.9)
  {a motor, and something for it to turn\\{\scriptsize and a mechanical fixture,
   which is usually the expensive part}};
\draw[hiz, ->] (nucleo.east) -- node[lbl, above]{six gate signals} (stage.west);
\draw[hiz, ->] (stage) -- (motor);

\node[umlnote, text width=54mm, anchor=north west] at (-4.6,-2.1)
  {\textbf{The rule that applies the moment a bridge is bought.} A half bridge
   whose two transistors conduct at the same time is a short circuit across the
   supply, and the dead time is the only thing preventing it. Measure the dead
   time with the loopback above, with no bridge connected, before the bridge is
   powered. Then power it with the supply current limited. Never adjust the dead
   time while the bridge is energised.};
\node[umlnote, text width=44mm, anchor=north west] at (3.4,-2.1)
  {The two loopback wires are the whole instrument. Capturing the timer's own
   output pins gives 3.57 ns resolution on a quantity that is usually measured
   with an oscilloscope, and this bench has none.};
\end{tikzpicture}
